\scanpage{5}
% L1–6
structure envisagée équivaut à la donnée de l'\uline{antipodisme} correspondant $\sigma=\sigma_H$ à cette structure $H$, qui \deleted{\unsure{doit}} \inserted{est} un autom \inserted{involutif} de la structure envisagée, induisant la transposition de $\unsure{\tau|T}$ i.e. échangeant entre eux les deux ens. $S_i$
% slanted
\margin{\deleted{en effet, tout élément d'ordre $2$ de \unsure{$G$}}}
% slanted
\margin{il faut exiger de plus que $\sigma$ opère trivialement sur le $\omega$ (défini comme le ``$\omega$ commun'' des $\omega_{S_i}$ $(i\in\tau)$) \unsure{$\Leftrightarrow$|$\rightleftarrows$} il y en a trois autres inv. qui induisent \deleted{\unsure{l'aut.}} la transp. antipodique de \unsure{$\tau$} \uline{et} de $\omega$, ils correspondent aux éléments $\unsure{t_\sigma|t_{\sigma'}}$ de $\bigwedge_{i\in\tau}S_i$ (produit contracté sous $\unsure{\Gamma_\omega|\mathscr T_\omega}=\mathbb Z/3\mathbb Z\wedge_{\pm1}\omega$)}
-- il y en a exactement $3$ correspondant (pour le choix d'un $i\in\tau$ \unsure{fixé}) aux trois \unsure{isom|isoms}
\[
S_i\simeq S_j\qquad(j\in\tau,\ j\ne i)
\]
qui ``conservent l'orientation''\unsure{)} i.e. induisent $\omega_{S_i}\simeq\omega_{S_j}$ l'iso. donné grâce à \unsure{$(*)$|$(A)$}. \deleted{Soit}

% L13–20
Alors l'ens $G_H$ des autom de $H=H_\sigma$ est le centralisateur de $\sigma$ -- $G_H$\inserted{$=Z(\sigma)$} est d'indice $3$ dans $G$, il est non distingué\inserted{ et en fait est son propre normalisateur}, ses trois conjugués correspondent aux trois conjugués de $\sigma$, qui sont les trois ``antipodismes'' de $G$, \unsure{correspondants|correspondant} aux trois structures hexagonales subordonnées à la structure ``trihexagonale'' envisagée.

% archive sheet number: 4
% L21–22
Considérons le sous-groupe (invariant) \inserted{$G'$} de $\deleted{\unsure{G'}}G$ engendré par \unsure{l'ens|l'un|l'ens.} $I$ des trois anti%
